%%                       Template LaTeX pour SSTIC
%% 
%% 
%% Utilisation
%% ===========
%% 
%% Pour compiler les sources, c'est aussi simple que :
%%     $ make
%% 
%% Le résultat est dans sstic-actes.pdf
%% 
%% Ce modèle doit être la base commune à tous les articles LaTeX pour
%% SSTIC.
%% 
%% Un exemple d'article est disponible dans le répertoire MonNom/ En
%% tant qu'auteur, vous devez copier ce répertoire et l'adapter.
%% 
%% Important : Vous ne devez pas pas toucher aux fichiers préfixés par _
%% 

\documentclass[b5paper,11pt]{llncs}

%% Nous avons déjà sélectionné quelques packages, si vous avez besoin
%% d'ajouter un nouveau package, veuillez contacter l'équipe (actes
%% (AT) lists.sstic.org) car nous devons vérifier la compatibilité des
%% modules entre eux.
%% 

\usepackage[french,british,UKenglish,english,USenglish,american]{babel}
\usepackage[utf8x]{inputenc}
\usepackage[fit]{truncate}
\usepackage{xspace}
\usepackage[table]{xcolor}
\usepackage{graphicx}
\usepackage{fancyhdr}
\usepackage{listings}
\usepackage{alltt}
\usepackage{ifthen}
\usepackage{url}

\usepackage{ean13isbn}

% Hack to have agenda working
\newcommand{\pause}[1]{\emph{#1}}

% Needed by a lot of people
\usepackage{tikz}
\usepackage{tikz-qtree}
\usetikzlibrary{arrows,positioning}
\usetikzlibrary{patterns}
\usetikzlibrary{shapes}


\usepackage{makeidx}
\makeindex

\usepackage[T1]{fontenc}
\usepackage{tabularx}
\usepackage{array}
\usepackage{multirow}
\usepackage{lscape}

\usepackage{subfigure}

\usepackage{etoolbox}
\usepackage[fancy]{tikz-inet}
\usetikzlibrary{shapes,backgrounds}

\usetikzlibrary{trees}

\usepackage{microtype}     % Needs the cm-super Debian package
\usepackage{lmodern}

\usepackage{verbatim}
\usepackage{amsmath}
\usepackage{amssymb}
\usepackage{pifont}
\usepackage{pmboxdraw}
\usepackage{fancyvrb}

\renewcommand\andname{et}
\renewcommand\lastandname{ et}

\newcommand{\xxx}{{\bf \LARGE \textcolor{red}{XXX}}\xspace}
\newcommand{\XXX}{\xxx}
\newtheorem{exx}{Exemple}

\def\keywordname{{\bf Mots-cl\'es:}}
\def\ackname{Remerciements.}
\def\andname{et}
\newcommand{\keywords}[1]{\par\addvspace\baselineskip\noindent\keywordname\enspace\ignorespaces#1}


\definecolor{grigri}{RGB}{233,232,234}
\lstset{breaklines=true,
        tabsize=4,
        boxpos=c,
        showstringspaces=false,
        basicstyle=\ttfamily\scriptsize,
        aboveskip=\bigskipamount,
        belowskip=\bigskipamount,
        captionpos=b,
        language=C,
        extendedchars=false
        basicstyle=\scriptsize\sffamily\color{darkgray},
        keywordstyle=\bfseries\color{purple},
        identifierstyle=\bfseries\color{black},
        commentstyle=\ttfamily\itshape\color{purple},
        stringstyle=\ttfamily\color{brown},
        showstringspaces=false,
        frame=leftline,
        backgroundcolor=\color{grigri},
        fillcolor=\color{grigri}}

\lstloadlanguages{[Sharp]C,[x86masm]Assembler,java,PHP,C,Python}

%%%%%%%%%%%%%%%%%%%%%%%%%%%%%%%%%%%%%%%%%%%%%%%%%%%%%%%%%%%%%%%%%%%%%%%%%%%%%%%%

% example d'usage:
%     \settitle[Mon titre court]{Mon titre super mégalong de la mort}
\newcommand{\theshorttitle}{}
\newcommand\settitle[2][]{%
 \title{#2}%
 \ifthenelse{\equal{#1}{}}%
  {\renewcommand{\theshorttitle}{#2}\toctitle{#2}\titlerunning{#2}}%
  {\renewcommand{\theshorttitle}{#1}\toctitle{#1}\titlerunning{#1}}%
}

\newcommand{\theshortauthor}{}
\newcommand\setauthor[2][]{%
 \author{#2}%
 \ifthenelse{\equal{#1}{}}%
  {\renewcommand{\theshortauthor}{#2}\tocauthor{#2}\authorrunning{#2}}%
  {\renewcommand{\theshortauthor}{#1}\tocauthor{#1}\authorrunning{#1}}%
}

\pagestyle{fancy}
\fancyhf{}
\setlength{\headheight}{16pt}
\fancyhead[LE,RO]{\thepage}
\fancyhead[RE]{\truncate{.90\headwidth}{\theshorttitle}}
\fancyhead[LO]{\theshortauthor}

\newcommand{\myclearpage}{\newpage{\pagestyle{empty}\cleardoublepage}}

%%%%%%%%%%%%%%%%%%%%%%%%%%%%%%%%%%%%%%%%%%%%%%%%%%%%%%%%%%%%%%%%%%%%%%%%%%%%%%%%

\bibliographystyle{plain}

\sloppy

\begin{document}
\selectlanguage{french}

\pagestyle{empty}
\frontmatter

\chapter*{Pr\'eface}


\lipsum[1-5]

\begin{flushright}
Bon symposium,

Gérard Menvussa, pour le comit{\'e} d'Organisation.
\end{flushright}

%%% Local Variables:
%%% mode: TeX-PDF
%%% TeX-master: "../_master"
%%% End:

\vfill

%%% Local Variables:
%%% mode: TeX-PDF
%%% TeX-master: "../_master"
%%% End:

\clearpage

\section*{Comit\'e d'organisation}
\begin{tabular}{@{}p{5cm}@{}p{6.5cm}@{}}
  Nicolas  \textsc{Bareil}        & Airbus              \\
  Mathieu  \textsc{Blanc}         & CEA/DAM             \\
  Pierre   \textsc{Capillon}      & ANSSI               \\
  Olivier  \textsc{Courtay}       & DGA-MI              \\
  Olivier  \textsc{Levillain}     & ANSSI               \\
  Benjamin \textsc{Morin}         & ANSSI               \\
  Nicolas  \textsc{Prigent}       & LSTI                \\
  Raphaël  \textsc{Rigo}          & Airbus              \\
  Sarah  \textsc{Zennou}          & Airbus              \\
  Frédéric  \textsc{Tronel}   & CentraleSupélec         \\
\end{tabular}

\section*{Comité de programme}
\begin{tabular}{@{}p{5cm}@{}p{6.5cm}@{}}
  Damien \textsc{Aumaitre}        & Quarkslab           \\
  Nicolas  \textsc{Bareil}        & Airbus              \\
  Mathieu \textsc{Blanc}          & CEA/DAM             \\
  Aurélien \textsc{Bordes}        & ANSSI               \\
  Jean-Marie \textsc{Borello}     & DGA-MI              \\
  Pierre   \textsc{Capillon}      & ANSSI               \\
  Olivier \textsc{Courtay}        & DGA-MI              \\
  Marion \textsc{Daubignard}      & ANSSI               \\
  Géraud \textsc{De Drouas}       & ANSSI               \\
  Fabrice \textsc{Desclaux}       & CEA/DAM             \\
  Isabelle \textsc{Kraemer}       & Orange              \\
  Colas \textsc{Le Guernic}       & DGA-MI              \\
  Olivier \textsc{Levillain}      & ANSSI               \\
  Thierry \textsc{Marinier}       & Lexfo               \\
  Xavier \textsc{Mehrenberger}    & Airbus              \\
  Benjamin \textsc{Morin}         & ANSSI               \\
  Camille \textsc{Mougey}         & CEA/DAM             \\
  Sarah \textsc{Nataf}            & Orange              \\
  Nicolas \textsc{Prigent}        & LSTI                \\
  Raphaël \textsc{Rigo}           & Airbus              \\
  Tiphaine \textsc{Romand-Latapie} & Airbus             \\
  Philippe \textsc{Teuwen}        & Quarkslab           \\
  Frédéric \textsc{Tronel}        & CentraleSupélec     \\
  Sarah \textsc{Zennou}           & Airbus              \\
\end{tabular}


%%% Local Variables:
%%% TeX-master: "../master"
%%% End:

\clearpage

\section*{Partenaires}

\begin{center}
ANSSI -  CEA - EADS - MISC/Le journal de la s{\'e}curit{\'e} informatique - OSSIR - Sup\'elec - Universit\'e de Rennes 1
\end{center}

\begin{figure}[h]
\begin{center}
\parbox{3cm}{\includegraphics[width=3cm]{_images/anssi}}
\hfill
%\hspace*{1.5cm}
\parbox{3cm}{\includegraphics[width=3cm]{_images/cea}}
\hfill
%\hspace*{1.5cm}
\parbox{3cm}{\includegraphics[width=3cm]{_images/eads}}
\end{center}
\vfill
\begin{center}
\parbox{3cm}{\includegraphics[width=3cm]{_images/technicolor}}
\hfill
\parbox{3cm}{\includegraphics[width=3cm]{_images/inria-rennes}}
\hfill
\parbox{3cm}{\includegraphics[width=3cm]{_images/misc}}
\end{center}
\vfill
\begin{center}
\parbox{3cm}{\includegraphics[width=3cm]{_images/uni-rennes1}}
\hfill
\parbox{3cm}{\includegraphics[width=3cm]{_images/supelec}}
\hfill
\parbox{3cm}{\includegraphics[width=3cm]{_images/ossir}}
\end{center}
\vfill
\begin{center}
\end{center}
\end{figure}

%%% Local Variables:
%%% mode: TeX-PDF
%%% TeX-master: "../master"
%%% End:

\myclearpage

%%% Local Variables:
%%% mode: TeX-PDF
%%% TeX-master: "../_master"
%%% End:


\tableofcontents
\mainmatter
\pagestyle{fancy}


\renewcommand{\thelstlisting}{\arabic{lstlisting}}
\part*{Conférences}
\addcontentsline{toc}{part}{Conférences}

%%%%%%%%%%%%%%%%%%%%%%%%%%%%%%%%%%%%%%%%%%%%%%%%%%%%%%%%%%%%%%%%%%%%%%%%%%%%%%%

% Default language is set to french
% Change to enable proper hyphenation of english language
% Do not remove or your language may be overridden by previous articles in the proceedings.
\selectlanguage{french} % possible values: french, british, UKenglish, USenglish, english, american

\settitle[Sécurité des trotinettes]{Sécurité des trotinettes : prise de contrôle des roues à distance}


% Pas besoin d'obfusquer les adresses mails, \email s'en chargera (XXX)
\setauthor[M.~MonNom, A.~MonAutreNom]{MonPrenom MonNom \and AutrePrenom AutreNom\\
  \email{monprenom.monnom@mycorp.com\\autreprenom.autrenom@mycorp.com}}

% XXX: Gérer le cas des compagnies multiples
\institute{MyCorp}

\maketitle
\index{MonNom, M.}
\index{MonAutreNom, A.}

\begin{abstract}
  Lorem ipsum dolor sit amet, consectetur adipisicing elit, sed do
  eiusmod tempor incididunt ut labore et dolore magna aliqua. Ut enim ad
  minim veniam, quis nostrud exercitation ullamco laboris nisi ut
  aliquip ex ea commodo consequat. Duis aute irure dolor in
  reprehenderit in voluptate velit esse cillum dolore eu fugiat nulla
  pariatur. Excepteur sint occaecat cupidatat non proident, sunt in
  culpa qui officia deserunt mollit anim id est laborum.
\end{abstract}

%% Découpage en sections
%% =====================
%% 
%% La classe llcns ne définit que 4 niveaux de sections :
%% 
%% * section
%% * subsection
%% * subsubsection
%% * paragraph
%% 

\section{Introduction}

\subsection{État de l'art}

Lorem ipsum dolor sit amet, consectetur adipisicing elit, sed do
eiusmod tempor incididunt ut labore et dolore magna aliqua. Ut enim ad
minim veniam, quis nostrud exercitation ullamco laboris nisi ut
aliquip ex ea commodo consequat. Duis aute irure dolor in
reprehenderit in voluptate velit esse cillum dolore eu fugiat nulla
pariatur. Excepteur sint occaecat cupidatat non proident, sunt in
culpa qui officia deserunt mollit anim id est laborum. Lorem ipsum
dolor sit amet, consectetur adipisicing elit, sed do eiusmod tempor
incididunt ut labore et \emph{dolore} magna aliqua. Ut enim ad minim
veniam, quis nostrud exercitation ullamco laboris nisi ut aliquip ex
ea commodo consequat. Duis aute irure dolor in reprehenderit in
voluptate velit esse cillum dolore eu fugiat nulla pariatur. Excepteur
sint occaecat cupidatat non proident, sunt in culpa qui officia
deserunt mollit anim id est laborum.

%% Images, figures
%% ===============
%% 
%% Les images doivent être placées dans un environnement figure de
%% façon à pouvoir y associer facilement une légende et un label.
%% 
%% Quelques notes:
%% 
%% * La légende doit être en dessous de l'image
%% 
%% * Le placement de l'image par LaTeX ne garantit pas que votre
%%   image sera exactement où vous souhaitiez. Remplacez donc les
%%   formules "l'image ci-dessous" mais plutôt par sa référence
%%   ("l'image~\ref{fig:monnom:archi}")
%% 
%% * Format de l'image
%%   - Compatible pdflatex : PNG, PDF
%%   - Les formats vectoriels (comme PDF) sont recommandés car il
%%     arrive de devoir redimensionner les images à l'édition
%%     finale
%% 
%% * Il est demandé de créér un répertoire img/ dédié à toutes les
%%   images
%% 
%% * La référence doit être préfixée par votre fig:monnom pour éviter
%%   les collisions entre auteurs
%% 
%% * Prenez soin de vos couleurs, SSTIC n'imprime qu'en dégradé de
%%   gris, votre choix de couleurs doit donc respecter cette
%%   contrainte.
%% 
%%   Votre article sera publié sur papier et sur notre site Web. Vous
%%   pouvez donc fournir les images dans les deux "formats". Nommez
%%   explicitement vos images avec le préfixe bw- (exemple: bw-archi.pdf)
%% 
%%   À défaut d'avoir une version papier, nous la convertirons à
%%   l'aide la commande suivante:
%%       $ convert archi.pdf -colorspace Gray bw-archi.pdf
%% 

\begin{figure}[ht]
  \centering
  \includegraphics[width=0.4\textwidth]{MonNom/img/archi}
  \caption{Légende de l'image}
  \label{fig:monnom:archi}
\end{figure}

Lorem ipsum dolor sit amet, consectetur adipisicing elit, sed do
eiusmod tempor incididunt ut labore et dolore magna aliqua. Ut enim ad
minim veniam, quis nostrud exercitation ullamco laboris nisi ut
aliquip ex ea commodo consequat. Duis aute irure dolor in
reprehenderit in voluptate velit esse cillum dolore eu fugiat nulla
pariatur. Excepteur sint occaecat cupidatat non proident, sunt in
culpa qui officia deserunt mollit anim id est laborum. Lorem ipsum
dolor sit amet, consectetur adipisicing elit, sed do eiusmod tempor
incididunt ut labore et dolore magna aliqua. Ut enim ad minim veniam,
quis nostrud exercitation ullamco laboris nisi ut aliquip ex ea
commodo consequat. Duis aute irure dolor in reprehenderit in voluptate
velit esse cillum dolore eu fugiat nulla pariatur. Excepteur sint
occaecat cupidatat non proident, sunt in culpa qui officia deserunt
mollit anim id est laborum.

%% Citations
%% =========
%% 
%% * Les références de citation doivent être préfixées par votre nom
%%   afin d'empêcher les collisions entre auteurs.
%% 
%% * La bibliographie doit être dans contenue dans le fichier
%%   biblio.bib au format BibTeX
%%   Référence des types : http://newton.ex.ac.uk/tex/pack/bibtex/btxdoc/node6.html
%%   Référence des champs: http://newton.ex.ac.uk/tex/pack/bibtex/btxdoc/node7.html
%% 
%% * Faîtes attention aux warning indiquant les références manquantes
%%
%% * Attention à la syntaxe lorsque vous citez plusieurs références:
%%   MAUVAIS: \cite{foo, bar}
%%   MAUVAIS: \cite{foo}\cite{bar}
%%   CORRECT: \cite{foo,bar}
%% 

Je sais aussi faire des citations, dans l'article de Charlie
Lembrouille, \cite{monnom:charlielembrouille}, il est démontré que, je cite,
\og{}les trotinettes sont vulnérables aux attaques par XSS\fg{}. C'est
moche, l'image~\ref{fig:monnom:archi} en étant la preuve !

Si une citation est manquante (comme \cite{monnom:referencefoireuse}),
un warning en rouge s'affichera lors de la compilation.

Avec des listes partout :

\begin{itemize}
\item Lorem ipsum dolor sit amet
  \begin{itemize}
  \item Parce que je le vaux bien
  \item N'est-ce pas ?
  \end{itemize}
\item Consectetur adipisicing elit
\item Sed do eiusmod tempor
\end{itemize}

Lorem ipsum dolor sit amet, consectetur adipisicing elit, sed do
eiusmod tempor incididunt ut labore et dolore magna aliqua. Ut enim ad
minim veniam, quis nostrud exercitation ullamco laboris nisi ut
aliquip ex ea commodo consequat. Duis aute irure dolor in
reprehenderit in voluptate velit esse cillum dolore eu fugiat nulla
pariatur. Excepteur sint occaecat cupidatat non proident, sunt in
culpa qui officia deserunt mollit anim id est laborum. Lorem ipsum
dolor sit amet, consectetur adipisicing elit, sed do eiusmod tempor
incididunt ut labore et dolore magna aliqua. Ut enim ad minim veniam,
quis nostrud exercitation ullamco laboris nisi ut aliquip ex ea
commodo consequat. Duis aute irure dolor in reprehenderit in voluptate
velit esse cillum dolore eu fugiat nulla pariatur. Excepteur sint
occaecat cupidatat non proident, sunt in culpa qui officia deserunt
mollit anim id est laborum.

\section{Use the source, Luke}

%% Texte verbatim
%% ==============
%% 
%% lstlisting
%% ----------
%% 
%% Documentation complète :  http://mirrors.ctan.org/macros/latex/contrib/listings/listings.pdf
%% 
%% Options utiles :
%% 
%% * ``language'' peut prendre les valeurs suivantes:
%%    Python	Java		PHP		[x86masm]Assembler
%%    C		Perl		HTML		[Sharp]C
%% 
%% * ``numbers'' pour numéroter les ligne (``stepnumbers'' pour  contrôler l'incrément)
%% 
%% * ``basicstyle'' pour modifier la police utilisée (exemple: basicstyle=\tiny)
%% 

Des fois, je peux aussi inclure du texte brut, comme la sortie
standard en figure~\ref{lst:monnom:outhello} qui a été généré par le
code~\ref{lst:monnom:helloword}.

\begin{lstlisting}[language={},caption={Sortie standard},label={lst:monnom:outhello}]
Hello world everybody
\end{lstlisting}

\begin{lstlisting}[language={Python},caption={Mon premier code},label={lst:monnom:helloword}]
#! /usr/bin/env python

import sys

def hello(name):
    print 'Hello world %s' % (name)

if __name__ == '__main__':
    hello(sys.argv[1] if len(sys.argv) > 1 else 'everybody')
\end{lstlisting}

\bibliography{MonNom/biblio}

%%% Local Variables:
%%% mode: TeX-PDF
%%% TeX-master: "master"
%%% End:

\myclearpage
\renewcommand{\thesection}{\arabic{section}}

%%%%%%%%%%%%%%%%%%%%%%%%%%%%%%%%%%%%%%%%%%%%%%%%%%%%%%%%%%%%%%%%%%%%%%%%%%%%%%%

%%% Local Variables:
%%% mode: TeX-PDF
%%% TeX-master: "master"
%%% End:


\myclearpage
\cleardoublepage
\newpage
\addtocmark[2]{Index des auteurs} % additional numbered TOC entry
\renewcommand{\indexname}{Index des auteurs}
\printindex

\end{document}

%% Publication
%% ===========
%% 
%% Vous devez envoyer les *sources* à actes (AT) lists.sstic.org avant
%% la date butoir.
%% 
%% L'auteur consciencieux prendre soin de corriger tous les warnings
%% avant de nous envoyer ses sources :)
%%
%% Vous pouvez envoyer une archive complète (.tar.gz), ou un patch, ou
%% un lien vers votre repository, etc.
%% 
%%% Local Variables:
%%% mode: TeX-PDF
%%% TeX-master: "master"
%%% End:
